\chapter{Introducere}
\label{ch:introducere}

\section{Context și Motivație}

Această lucrare este un studiu experimental în învățarea automată, mai precis în subdomeniul \emph{detecției de anomalii} (\emph{anomaly detection}) -- sarcina de a detecta puncte, evenimente sau structuri care deviază de la normalitate. Domeniul este unul larg și acoperă o varietate de probleme, dintre care în această lucrare sunt abordate: detecția fraudelor în tranzacții bancare, detecția atacurilor asupra rețelelor de calculatoare și detecția defecțiunilor în metricile monitorizate ale serverelor.

O primă distincție este cea dintre \emph{outlier detection}, care antrenează pe date ce conțin deja anomalii, și \emph{novelty detection}, care antrenează exclusiv pe date normale și evaluează, la testare, capacitatea modelului de a detecta devieri de la normalitatea astfel învățată. Această lucrare adoptă abordarea \emph{novelty detection}, cadrul folosit practic în cazurile de utilizare de mai sus (un model de detecție a fraudelor nu este antrenat pe cazuri confirmate de fraudă), care ajută și la formularea scopului lucrării printr-o singură întrebare: în ce măsură modul în care un model reprezintă structura latentă a datelor \enquote{normale} determină tipurile de anomalii pe care le poate detecta?

Motivația, cât și scopul acestei lucrări sunt, așadar, căutarea unui răspuns la această întrebare și abordarea ulterioară a ipotezelor formate în urma rezultatelor experimentale.

Contribuția proprie a acestei lucrări constă, așadar, în:

\begin{enumerate}
    \item un pipeline standardizat de preprocesare, antrenare și evaluare, aplicat uniform pe trei seturi de date diferite și pe până la treisprezece configurații de modele și seturi de caracteristici (\emph{features});
    \item o analiză sistematică a performanței diferitelor paradigme în funcție de tipul setului de date; și
    \item o serie de studii de caz de explicabilitate, menite să testeze direct ipoteze formulate în secțiunea de discuții a rezultatelor, folosind hărți termice ale erorii de reconstrucție, o măsură de concentrație bazată pe coeficientul Gini și valori de atribuire SHAP, în locul unor concluzii bazate exclusiv pe metrici agregate.
\end{enumerate}

\section{Structura Lucrării}

Restul lucrării este organizat în cinci capitole: 
Capitolul~\ref{ch:preliminarii} (\emph{Preliminarii}) introduce fundalul teoretic -- funcționarea și paradigma fiecărui model testat, și seturile de date folosite; 
Capitolul~\ref{ch:methodology} (\emph{Metode și Configurare Experimentală}) descrie pipeline-urile de preprocesare, alegerile de hiperparametri/arhitectură și protocolul de evaluare; 
Capitolul~\ref{ch:results} (\emph{Rezultate Experimentale}) prezintă rezultatele pe seturi de date, cu discuții care le raportează la așteptările din Preliminarii; 
Capitolul~\ref{ch:explainability} (\emph{Studii de Caz de Explicabilitate}) revizitează și testează direct, prin diagnostice țintite, ipoteze specifice ridicate în discuția Capitolului~\ref{ch:results}, confirmând unele și nuanțând altele; 
iar Capitolul~\ref{ch:concluzii} (\emph{Concluzii}) rezumă concluziile, discută limitările și schițează direcții de dezvoltare ulterioară.